\documentclass[12pt]{article}

\usepackage{fontspec}
% надёжный шрифт с кириллицей (работает в Overleaf)
\setmainfont{Libertinus Serif}[Ligatures=TeX]
\setsansfont{Libertinus Sans}[Ligatures=TeX]
\newfontfamily\cyrillicfont{Libertinus Serif}
\newfontfamily\cyrillicfontsf{Libertinus Sans}
\usepackage{polyglossia}
\setmainlanguage{russian}
\setotherlanguage{english}

\usepackage{amsmath}
\usepackage{amsfonts}
\usepackage{amssymb}
\usepackage{graphicx}
\usepackage{microtype}

\title{Подробный вывод формул оптимального импульсного управления \\ для transfers между соседними орбитами}
\author{На основе работы T. N. Edelbaum}
\date{}

\begin{document}

\maketitle

\section{Постановка задачи}

Рассмотрим задачу оптимального импульсного перехода между двумя соседними орбитами с малой эксцентричностью. Ориентации линий апсид и узлов произвольны. Используем линеаризованные уравнения движения относительно круговой орбиты.

\section{Уравнения движения}

Введем переменные состояния:
\begin{align*}
x_1 &= \frac{da}{a_0} \\
x_2 &= e_y \\
x_3 &= e_x \\
x_4 &= i_y \\
x_5 &= i_x
\end{align*}

Линеаризованные уравнения изменения параметров:
\begin{align}
\frac{dx_1}{du} &= 2\sqrt{\frac{a_0}{K}}l_T \tag{1} \\
\frac{dx_2}{du} &= \sqrt{\frac{a_0}{K}}(2l_T\sin\theta - l_R\cos\theta) \tag{2} \\
\frac{dx_3}{du} &= \sqrt{\frac{a_0}{K}}(2l_T\cos\theta + l_R\sin\theta) \tag{3} \\
\frac{dx_4}{du} &= \sqrt{\frac{a_0}{K}}l_N\sin\theta \tag{4} \\
\frac{dx_5}{du} &= \sqrt{\frac{a_0}{K}}l_N\cos\theta \tag{5}
\end{align}

где $u$ - интеграл от ускорения тяги по времени, $\theta$ - центральный угол.

\section{Формулировка вариационной задачи}

Гамильтониан системы:
\begin{equation}
H = \sum_{i=1}^5 \lambda_i \frac{dx_i}{du} \tag{6}
\end{equation}

Сопряженные переменные $\lambda_i$ удовлетворяют уравнениям:
\begin{equation}
\frac{d\lambda_i}{du} = -\frac{\partial H}{\partial x_i}
\end{equation}

Введем вектор праймера с компонентами:
\begin{align}
\lambda &= \sqrt{\frac{a_0}{K}}(-\lambda_2\cos\theta + \lambda_3\sin\theta) \tag{8} \\
\mu &= \sqrt{\frac{a_0}{K}}(2\lambda_1 + 2\lambda_2\sin\theta + 2\lambda_3\cos\theta) \tag{9} \\
\nu &= \sqrt{\frac{a_0}{K}}(\lambda_4\sin\theta + \lambda_5\cos\theta) \tag{10}
\end{align}

Оптимальное управление:
\begin{equation}
l_R = \frac{\lambda}{p}, \quad l_T = \frac{\mu}{p}, \quad l_N = \frac{\nu}{p} \tag{7}
\end{equation}
где $p = \sqrt{\lambda^2 + \mu^2 + \nu^2}$.

\section{Преобразование уравнений для праймера}

Введем угол $\theta_0$:
\begin{equation}
\tan\theta_0 = \frac{\lambda_2}{\lambda_3} \tag{14}
\end{equation}

Тогда:
\begin{align}
\lambda &= \sqrt{\frac{a_0}{K}}\sqrt{\lambda_2^2 + \lambda_3^2}\sin(\theta - \theta_0) \tag{11} \\
\mu &= \sqrt{\frac{a_0}{K}}\left[2\lambda_1 + 2\sqrt{\lambda_2^2 + \lambda_3^2}\cos(\theta - \theta_0)\right] \tag{12}
\end{align}

Для нормальной компоненты:
\begin{align}
\nu &= \sqrt{\frac{a_0}{K}}\left[\frac{\lambda_3\lambda_4 - \lambda_2\lambda_5}{\sqrt{\lambda_2^2 + \lambda_3^2}}\sin(\theta - \theta_0) + \frac{\lambda_2\lambda_4 + \lambda_3\lambda_5}{\sqrt{\lambda_2^2 + \lambda_3^2}}\cos(\theta - \theta_0)\right] \tag{13}
\end{align}

\section{Три класса решений}

\subsection{Узловой случай (Nodal Case)}

Характеризуется условиями:
\begin{align}
\lambda_1 &= 0 \tag{15} \\
\theta_2 &= \theta_1 + \pi \tag{16} \\
\lambda(\theta_1) &= -\lambda(\theta_2) \\
\mu(\theta_1) &= -\mu(\theta_2) \\
\nu(\theta_1) &= -\nu(\theta_2) \tag{17}
\end{align}

Изменения орбитальных элементов:
\begin{align}
\frac{\Delta x_1}{u} &= 2\sqrt{\frac{a_0}{K}}l_{T_1}(1-2\delta) \tag{32} \\
\frac{\Delta x_2}{u} &= -\sqrt{\frac{a_0}{K}}l_{R_1} \tag{33} \\
\frac{\Delta x_3}{u} &= 2\sqrt{\frac{a_0}{K}}l_{T_1} \tag{34} \\
\frac{\Delta x_4}{u} &= 0 \tag{35} \\
\frac{\Delta x_5}{u} &= \sqrt{\frac{a_0}{K}}l_{N_1} \tag{36}
\end{align}

Суммарный импульс:
\begin{equation}
u = \sqrt{\frac{K}{a_0}}\sqrt{\Delta i^2 + \frac{\Delta e_x^2}{4} + \Delta e_y^2} \tag{37}
\end{equation}

\subsection{Невырожденный случай (Nondegenerate Case)}

Условия:
\begin{align}
\lambda_2\lambda_4 &= -\lambda_3\lambda_5 \tag{18} \\
\theta_2 - \theta_0 &= \theta_0 - \theta_1 \tag{19} \\
\lambda(\theta_1) &= -\lambda(\theta_2) \\
\mu(\theta_1) &= \mu(\theta_2) \\
\nu(\theta_1) &= -\nu(\theta_2) \tag{20}
\end{align}

Введем новые переменные:
\begin{align}
R &\equiv \sqrt{\frac{a_0}{K}}\cdot\sqrt{\lambda_2^2 + \lambda_3^2} \tag{50} \\
C &\equiv \cos(\theta - \theta_0) \\
S &\equiv \sin(\theta - \theta_0)
\end{align}

Направляющие косинусы импульсов:
\begin{align}
l_R &= \frac{2R^2S}{\sqrt{1+R^2}\sqrt{4R^2 + (1-3R^2)C^2}} \tag{52} \\
l_T &= \frac{\sqrt{1+R^2}C}{\sqrt{4R^2 + (1-3R^2)C^2}} \tag{53} \\
l_N &= \frac{2RS}{\sqrt{1+R^2}\sqrt{4R^2 + (1-3R^2)C^2}} \tag{54}
\end{align}

\subsection{Сингулярный случай (Singular Case)}

Условия:
\begin{align}
\lambda_1 &= 0 \tag{24} \\
\lambda_2\lambda_4 &= -\lambda_3\lambda_5 \tag{25} \\
3(\lambda_2^2 + \lambda_3^2)^2 &= (\lambda_3\lambda_4 - \lambda_2\lambda_5)^2 = \frac{3}{4}\frac{K}{a_0} \tag{26}
\end{align}

Компоненты праймера:
\begin{align}
\lambda &= \frac{1}{2}\sin(\theta - \theta_0) \tag{27} \\
\mu &= \cos(\theta - \theta_0) \tag{28} \\
\nu &= \frac{\sqrt{3}}{2}\sin(\theta - \theta_0) \tag{29}
\end{align}

\section{Вывод формул для импульсов}

\subsection{Узловой случай}

Из уравнений (32)-(36) выражаем компоненты импульса через изменения элементов:

\begin{align*}
l_{T_1} &= \frac{\Delta x_3}{2u}\sqrt{\frac{K}{a_0}} \\
l_{R_1} &= -\frac{\Delta x_2}{u}\sqrt{\frac{K}{a_0}} \\
l_{N_1} &= \frac{\Delta x_5}{u}\sqrt{\frac{K}{a_0}}
\end{align*}

Учитывая, что $l_{R_1}^2 + l_{T_1}^2 + l_{N_1}^2 = 1$, получаем:
\begin{equation*}
\left(\frac{\Delta x_2}{u}\right)^2\frac{K}{a_0} + \left(\frac{\Delta x_3}{2u}\right)^2\frac{K}{a_0} + \left(\frac{\Delta x_5}{u}\right)^2\frac{K}{a_0} = 1
\end{equation*}

Отсюда:
\begin{equation*}
u^2 = \frac{K}{a_0}\left(\Delta x_2^2 + \frac{\Delta x_3^2}{4} + \Delta x_5^2\right)
\end{equation*}

что соответствует уравнению (37).

\subsection{Невырожденный случай}

После введения переменных $R$, $C$, $S$ и решения системы уравнений получаем выражение для суммарного импульса:
\begin{equation}
u = \sqrt{\frac{K}{2a_0}}\sqrt{\Delta i^2 + \Delta e_x^2 + \Delta e_y^2 - \frac{\Delta a^2}{2a_0^2} + \sqrt{\left(\Delta i^2 - \Delta e_x^2 - \Delta e_y^2 + \frac{\Delta a^2}{a_0^2}\right)^2 + 4\Delta i^2\Delta e_y^2}} \tag{68}
\end{equation}

\subsection{Сингулярный случай}

Для сингулярного случая получаем:
\begin{equation}
u = \sqrt{\frac{K}{4a_0}}\sqrt{\left(\sqrt{3}\Delta i + \Delta e_y\right)^2 + \Delta e_x^2} \tag{75}
\end{equation}

\section{Заключение}

В работе представлен подробный вывод формул оптимального импульсного управления для переходов между соседними орбитами. Рассмотрены три класса решений: узловой, невырожденный и сингулярный. Для каждого случая получены выражения для суммарного импульса и оптимальных направлений thrust.

\end{document}